\documentclass[11pt]{article}

\usepackage[margin=1in]{geometry}
\usepackage{amsmath, amssymb}
\usepackage{xcolor}
\usepackage{listings}
\usepackage{hyperref}
\usepackage{array}
\usepackage{booktabs}

\hypersetup{
    colorlinks=true,
    linkcolor=blue!60!black,
    urlcolor=blue!60!black,
    citecolor=blue!60!black
}

\lstdefinestyle{bugpython}{
  language=Python,
  basicstyle=\ttfamily\small,
  keywordstyle=\color{blue!60!black},
  stringstyle=\color{green!40!black},
  commentstyle=\color{gray!70!black},
  showstringspaces=false,
  breaklines=true,
  frame=single,
  rulecolor=\color{gray!30},
  columns=fullflexible,
  keepspaces=true
}

\title{SymPy Bug Report}
\date{}

\begin{document}

\author{}
\maketitle

\section{Bug 52: \texttt{function\_range} Says \texttt{sqrt(tan(x)**2)} Is Constant}

\subsection{Minimal Reproducer}

\medskip

\begin{lstlisting}[style=bugpython]
from sympy import Interval, Rational, N, sqrt, tan, symbols
from sympy.calculus.util import function_range

x = symbols("x")
expr = sqrt(tan(x)**2)
actual = function_range(expr, x, Interval(-1, 1))

print("function_range =", actual)
print("value_at_0 =", expr.subs(x, 0))
print("value_at_1_over_2 =", N(expr.subs(x, Rational(1, 2)), 50))
print("range_contains_0 =", actual.contains(0))
print("expected_range =", Interval(0, tan(1)))
\end{lstlisting}

\subsection{Observed Behavior}

\medskip

\begin{lstlisting}[style=bugpython]
function_range = {tan(1)}
value_at_0 = 0
value_at_1_over_2 = 0.54630248984379051325517946578028538329755172017979
range_contains_0 = False
expected_range = Interval(0, tan(1))
\end{lstlisting}

SymPy returns the singleton set \(\{\tan(1)\}\). This says the expression is
constant on \([-1,1]\), but the same expression evaluates to \(0\) at \(x=0\)
and to a positive interior value at \(x=1/2\). Thus the returned set omits real
values attained by the input expression.

\subsection{Expected Behavior}

The mathematically correct result is
\[
  \operatorname{function\_range}\!\left(\sqrt{\tan(x)^2}, x, [-1,1]\right)
  = [0,\tan(1)].
\]
In SymPy notation, the expected output is \texttt{Interval(0, tan(1))}.

\subsection{Mathematical Explanation}

The interval \([-1,1]\) lies strictly inside the principal continuity interval
\((-\pi/2,\pi/2)\) of \(\tan(x)\). On this interval, \(\tan(x)\) is continuous
and strictly increasing, with \(\tan(0)=0\), \(\tan(-1)<0\), and \(\tan(1)>0\).
For real \(x\),
\[
  \sqrt{\tan(x)^2}=|\tan(x)|.
\]
Therefore the minimum is attained at the interior point \(x=0\), where the value
is \(0\), and the maximum is attained at the endpoints, where
\[
  |\tan(-1)| = \tan(1), \qquad |\tan(1)| = \tan(1).
\]
The range is consequently the full interval \([0,\tan(1)]\). The singleton
\(\{\tan(1)\}\) is not an equivalent representation, because it excludes the
attained value \(0\) and every interior value such as \(|\tan(1/2)|\).

\subsection{Source-Level Diagnosis}

The diagnosis identifies the relevant implementation as
\nolinkurl{sympy/calculus/util.py}, in the public \texttt{function\_range}
routine. The traced call path enters \texttt{function\_range} around line 161,
computes the continuous domain with \texttt{continuous\_domain} around line 233,
then in the interval branch collects endpoint values around lines 253--258 and
critical points from
\[
  \texttt{solveset(f.diff(symbol), symbol, interval)}
\]
around line 260. The final range is then constructed from the infimum and
supremum of the collected values around line 281.

For this input, \texttt{continuous\_domain} returns \([-1,1]\), so the
expression itself is continuous on the requested interval. However, its
derivative is
\[
  \frac{(2\tan(x)^2+2)\sqrt{\tan(x)^2}}{2\tan(x)},
\]
which is undefined at \(x=0\). The derivative-zero solve over \([-1,1]\)
returns the empty set, and the implementation does not also include finite
points where the derivative is undefined while the original expression remains
defined and continuous. As a result, the only sampled values are the two endpoint
values \(f(-1)=\tan(1)\) and \(f(1)=\tan(1)\). The interior cusp value
\(f(0)=0\) is never sampled, so the range degenerates to \(\{\tan(1)\}\).

A plausible fix direction supported by the diagnosis is to include finite
interior points where \texttt{f.diff(symbol)} is undefined but \texttt{f} is
defined and continuous. For this example, the zero of the derivative denominator
\(\tan(x)\) would add \(x=0\) as a critical point.

\subsection{Additional Failing Instances}

The same failure occurs for symmetric intervals \([-b,b]\) contained in
\((-\pi/2,\pi/2)\). In each case,
\[
  \sqrt{\tan(x)^2}=|\tan(x)|,
\]
so the minimum remains \(0\) at the cusp \(x=0\), while the endpoint values are
both \(\tan(b)\). SymPy instead returns only the endpoint value because it misses
the derivative-undefined but continuous critical point at the origin.

\begin{center}
\small
\renewcommand{\arraystretch}{1.3}
\begin{tabular}{@{}>{\raggedright\arraybackslash}p{0.44\linewidth}>{\raggedright\arraybackslash}p{0.23\linewidth}>{\raggedright\arraybackslash}p{0.23\linewidth}@{}}
\toprule
\textbf{Input / Expression} & \textbf{SymPy behavior} & \textbf{Expected behavior} \\
\midrule
\(\sqrt{\tan(x)^2}\) on \([-1,1]\) & \(\{\tan(1)\}\) & \([0,\tan(1)]\) \\
\(\sqrt{\tan(x)^2}\) on \([-1/2,1/2]\) & \(\{\tan(1/2)\}\) & \([0,\tan(1/2)]\) \\
\(\sqrt{\tan(x)^2}\) on \([-1/3,1/3]\) & \(\{\tan(1/3)\}\) & \([0,\tan(1/3)]\) \\
\(\sqrt{\tan(x)^2}\) on \([-2/3,2/3]\) & \(\{\tan(2/3)\}\) & \([0,\tan(2/3)]\) \\
\(\sqrt{\tan(x)^2}\) on \([-3/4,3/4]\) & \(\{\tan(3/4)\}\) & \([0,\tan(3/4)]\) \\
\(\sqrt{\tan(x)^2}\) on \([-1/4,1/4]\) & \(\{\tan(1/4)\}\) & \([0,\tan(1/4)]\) \\
\bottomrule
\end{tabular}
\end{center}

\subsection{Related Issues Assessment}

The best-effort deduplication check found related public reports involving
\texttt{function\_range}, absolute values, trigonometric expressions, and
fragile critical-point handling. The closest matches include a \texttt{TypeError}
tied to a \texttt{ConditionSet} for an expression involving
\(\operatorname{Abs}\), a separate inverse-trigonometric sign or singularity
issue, and a slow \texttt{function\_range} case for a tangent expression. The
available evidence did not identify a public report for \(\sqrt{\tan(x)^2}\) on
\([-1,1]\), the wrong singleton output \(\{\tan(1)\}\), or the missed continuous
cusp value \(0\). The verdict is therefore a new instance of a known failure
mode: the general fragility is related to known behavior, but this exact
reproducible case is previously unreported and remains individually actionable
as an issue and regression test.

\subsection{Proposed SymPy Regression Test}

\medskip

\begin{lstlisting}[style=bugpython]
import pytest

from sympy import Interval, Rational, S, sqrt, tan, symbols
from sympy.calculus.util import function_range


@pytest.mark.parametrize(
    "bound",
    [S.One, Rational(1, 2), Rational(1, 3), Rational(2, 3), Rational(3, 4), Rational(1, 4)],
)
def test_function_range_sqrt_tan_squared_includes_cusp_value(bound):
    x = symbols("x")
    assert function_range(sqrt(tan(x)**2), x, Interval(-bound, bound)) == Interval(0, tan(bound))
\end{lstlisting}

\end{document}
